% Generated draft: requires mathematical and source review.
\documentclass[11pt]{article}
\usepackage[margin=25mm]{geometry}
\usepackage{fontspec}
\setmainfont{Latin Modern Roman}
\usepackage{amsmath,amssymb,mathtools,unicode-math}
\setmathfont{Latin Modern Math}
\usepackage{enumitem,needspace,xurl,hyperref}
\setlength{\parindent}{0pt}
\setlength{\parskip}{5pt}
\setlength{\emergencystretch}{3em}
\newcommand{\N}{\mathbb{N}}
\newcommand{\Z}{\mathbb{Z}}
\newcommand{\Q}{\mathbb{Q}}
\newcommand{\R}{\mathbb{R}}
\newcommand{\C}{\mathbb{C}}
\newcommand{\F}{\mathbb{F}}
\newcommand{\A}{\mathbb{A}}
\newcommand{\PP}{\mathbb{P}}
\newcommand{\E}{\mathbb{E}}
\newcommand{\eps}{\varepsilon}
\newcommand{\dd}{\,\mathrm{d}}
\begin{document}
\section*{UnsolvedMath ID: 800935}
\textbf{Source problem code:} catalogue entry
\subsection{Definitions and mathematical statement}
Let \(M\) be a smooth closed manifold, where closed means compact without boundary. Let \(\mathcal{C}(M)\) denote the set of free homotopy classes of loops in \(M\). For a Riemannian metric \(g\), define its marked length spectrum by
\[
\ell_g\colon\mathcal{C}(M)\longrightarrow[0,\infty),\qquad
\ell_g(c)=\inf_{\substack{\gamma\colon S^1\to M\text{ piecewise smooth}\\[\gamma]=c}}\operatorname{Length}_g(\gamma).
\]
The marking here is the identification of classes on the same underlying manifold \(M\).

Interpreting the source's phrase ``negative curvature'' as strictly negative sectional curvature, the proposed statement is: for every such \(M\) and every pair of Riemannian metrics \(g_0,g_1\) satisfying
\[
K_{g_i}(\sigma)<0\qquad
(i\in\{0,1\},\ x\in M,\ \sigma\subset T_xM\text{ a two-dimensional subspace}),
\]
one has
\[
\bigl(\forall c\in\mathcal{C}(M),\ \ell_{g_0}(c)=\ell_{g_1}(c)\bigr)
\quad\Longrightarrow\quad
\bigl(\exists F\colon M\to M\text{ a diffeomorphism},\ F^*g_1=g_0\bigr).
\]
The source asks only for an isometry; it does not explicitly require \(F\) to be homotopic to the identity.

The infimum definition is used throughout. In particular, the class of a constant loop has value zero. The source's parenthetical description as the length of a shortest closed geodesic requires a convention allowing constant geodesics in this class, or a restriction of that description to nontrivial classes. The source initially assumes that \(M\) is closed and later says compact; the formulation above retains the initial absence of boundary. The supplied registry labels this conjecture open; this is a supplied status, not an independently verified mathematical assessment.
\subsection{Short English statement}
On a closed manifold, two metrics of strictly negative sectional curvature are conjectured to be isometric if the infimum of loop lengths agrees in every free homotopy class. The source does not specify a homotopy requirement on the isometry. This is a statement-only draft requiring human mathematical review.
\subsection{Sources}
\begin{itemize}
\item Local UnsolvedMath registry, supplied record with permanent ID \(800935\), exact source problem code \texttt{catalogue entry}. The record lists Burns and Matveev as proposers, gives the year 2019, and labels the problem open.
\item Burns, Matveev, \emph{Open Problems and Questions about Geodesics} (2019), Conjecture 15, as identified in the supplied record: \url{https://arxiv.org/abs/1308.5417}. The supplied provenance records access on 2026-10-16 and extraction by \texttt{arxiv\_latex}; these are supplied provenance details, not a claim of independent source access or verification.
\item The supplied title and background give the attribution/context label \texttt{[Bur-Kat1985]}. No full bibliographic entry for this label is supplied, so its authorship and publication details remain unresolved.
\end{itemize}
\end{document}
